% Generated draft: requires mathematical and source review.
\documentclass[11pt]{article}
\usepackage[margin=25mm]{geometry}
\usepackage{fontspec}
\setmainfont{Latin Modern Roman}
\usepackage{amsmath,amssymb,mathtools,unicode-math}
\setmathfont{Latin Modern Math}
\usepackage{enumitem,needspace,xurl,hyperref}
\setlength{\parindent}{0pt}
\setlength{\parskip}{5pt}
\setlength{\emergencystretch}{3em}
\newcommand{\N}{\mathbb{N}}
\newcommand{\Z}{\mathbb{Z}}
\newcommand{\Q}{\mathbb{Q}}
\newcommand{\R}{\mathbb{R}}
\newcommand{\C}{\mathbb{C}}
\newcommand{\F}{\mathbb{F}}
\newcommand{\A}{\mathbb{A}}
\newcommand{\PP}{\mathbb{P}}
\newcommand{\E}{\mathbb{E}}
\newcommand{\eps}{\varepsilon}
\newcommand{\dd}{\,\mathrm{d}}
\begin{document}
\section*{UnsolvedMath ID: 11000036}
\textbf{Source problem code:} catalogue entry
\subsection{Definitions and mathematical statement}
Let \(G=\operatorname{Mod}_g\) and \(P\subseteq G\) be the pseudo-Anosov mapping classes, those admitting a representative with transverse invariant measured foliations expanded and contracted by factors \(\lambda\) and \(\lambda^{-1}\), \(\lambda>1\). The printed conjecture is
\[
 d(P)=1,
\]
where the preceding definition of \(d\) is missing.

If the intended density is uniform word-ball density for a fixed finite generating set \(S\), define \(B_S(R)=\{h:\|h\|_S\leq R\}\). In that interpretation the precise assertion is existence of the limit and
\[
 \lim_{R\to\infty}\frac{|P\cap B_S(R)|}{|B_S(R)|}=1.
\]
This is an explicitly conditional interpretation: the source does not repeat the exhaustion, limit convention or generating-set quantifier.

The accompanying source distinguishes this conjecture from J. Maher's result that a suitable random walk lands in \(P\) with probability tending to \(1\), and I. Rivin's result for a particular sampling scheme via \(\operatorname{Sp}(2g,\Z)\). It expressly says that the conjecture does not appear to follow directly from those results. A statement \(\mu^{*n}(P)\to1\) for a walk law is not automatically the displayed uniform-ball statement. The supplied research summary calls the problem solved in a generic/density sense conditional on its interpretation, but supplies no exact theorem identifying that sense with \(d\). That discrepancy remains unresolved rather than being used to declare the printed conjecture solved.
\subsection{Short English statement}
Pseudo-Anosov mapping classes should have density one in the source's intended sense. Random-walk genericity and the possible uniform word-ball formulation must be distinguished.
\subsection{Sources}
\begin{itemize}
\item \emph{Problems on Mapping Class Groups and Related Topics}, edited by Farb, Conjecture 3.15, ``Density of pseudo-Anosovs,'' supplied PDF page 36. \url{https://www.math.uchicago.edu/~farb/papers/mcgbook.pdf}. Recorded access 2026-07-29; extraction pdf-font-structure.
\item J. Maher [Mah] and I. Rivin [Ri], source references for the random-walk and specific symplectic-sampling results. Full bibliography entries and exact hypotheses are absent. The summary mentions 2010 and later results without a precise uniform-word-ball theorem citation.
\item Local UnsolvedMath registry, \texttt{data/unsolvedmath/problems.json}, permanent ID \(11000036\), exact source code \texttt{catalogue entry}; research classification PARTIAL-PROGRESS, with a conditional generic-sense solved report in its summary.
\end{itemize}
\end{document}
